\documentclass[11pt]{article}

\usepackage{series}   % house style
\usepackage{amsmath,amssymb,amsfonts}
\usepackage{geometry}
\usepackage{hyperref}
\usepackage{enumitem}


\geometry{margin=1in}

\theoremstyle{plain}
\newtheorem{axiom}{Axiom}[section]
\theoremstyle{plain}
\newtheorem{postulate}{Postulate}[section]


% ---------------- Metadata ----------------
\title{The Modal Triplet Theory Program A0:\\ A Structural Theory of Reduced Description}
\author{Peter Nero}
\date{July 2026}

\begin{document}

\maketitle

\begin{abstract}
We give a typed, conditional formulation of Modal Triplet Theory (MTT) as a
theory of reduced description.  A representative section, an exact upper
decoder, autonomous reduced evolution, and merger of effective states solve
four different mathematical problems; none follows merely from
noninjectivity of a projection.  We prove the factor-through criterion for
autonomous descent, a finite-diameter obstruction for a genuinely
contractive reduced self-map, exact basin-local fixed-point results, and
controlled approximate-orbit estimates.  Admissibility boundaries are
defined by explicit margins and section or projector conditioning rather
than by assuming the desired obstruction.  The encoding atlas and its
relations are consequently local and conditional.  Probability,
irreversibility, physical time, and any concrete physical realization
require their own measure, evolution, and source hypotheses.
\end{abstract}

\section*{Revision note for version 2}
\addcontentsline{toc}{section}{Revision note for version 2}
\paragraph{Supersedes.}
Version~1 of Program A0.
\paragraph{Reason.}
The former text conflated representative selection with
microscopic recovery, treated an additive-error estimate as a Banach
contraction, and promoted several conditional atlas statements to universal
theorems.
\paragraph{Resolution.}
Version~2 installs the projection--descent and recovery
classification, exact basin-local FCC, the finite-diameter self-map
obstruction, explicit margin barriers, and a typed distinction between
physical evolution and auxiliary stabilization.  It also corrects the
direction of encoding factorization and scopes probability, records,
selection, and global relations.
\paragraph{Retained content.}
Finite admissibility, local encoding charts,
coherence margins, and the total chart-of-charts remain the structural core.
\paragraph{Open boundary.}
No physical measure, Born rule, arrow of time, reset
law, spacetime dynamics, or concrete MTT realization is derived in A0.

\tableofcontents
\newpage

% =====================================================
\section*{Part I: Abstract Framework}
\addcontentsline{toc}{section}{Part I: Abstract Framework}

\section{Typed Data and Admissibility}

This section fixes notation and logical order.  All definitions are
mathematical; physical interpretation is an additional assignment.

\subsection{Upper evolution and reductions}

\begin{definition}[Typed projection system]
A typed projection system at parameter $t$ consists of
\[
(X_0,X_t,\Phi_t;\;Y_0,Y_t,P_0,P_t)
\]
where:
\begin{itemize}
\item $X_0$ and $X_t$ are standard Borel spaces, equipped with metrics when
metric estimates are used;
\item $\Phi_t:X_0\to X_t$ is a measurable upper evolution map;
\item $P_0:X_0\to Y_0$ and $P_t:X_t\to Y_t$ are measurable reductions,
surjective onto their declared effective images.
\end{itemize}
The cross-level output map is
\[
Q_t:=P_t\circ\Phi_t:X_0\longrightarrow Y_t.
\]
\end{definition}

\begin{remark}[Evolution semantics]
No invertibility is needed for the typing results below.  In an application,
$\Phi_t$ must be declared either as physical evolution, with its own
well-posedness, causal or hyperbolic hypotheses, or as an auxiliary
stabilization map or semigroup.  An auxiliary parameter is not physical time,
and contraction of an auxiliary semigroup is not by itself physical
irreversibility.
\end{remark}

\subsection{Tolerance}

\begin{definition}[Tolerance]
Fix $\varepsilon>0$.
Two points $y_1,y_2$ in a declared reduced metric space are
$\varepsilon$-indistinguishable if
\[
d_Y(y_1,y_2)\le \varepsilon,
\]
where $d_Y$ is a chosen pseudometric.
\end{definition}

\subsection{Representative selection, recovery, and descent}

\begin{definition}[Representative section]
A representative section for $Q_t$ is a map $S_t:Y_t\to X_0$ satisfying
\[
Q_t\circ S_t=\operatorname{id}_{Y_t}.
\]
It chooses one compatible upper representative.  It does not recover the
upper state that actually produced the reduced datum.
\end{definition}

\begin{definition}[Exact upper decoder]
An exact upper decoder is a map $D_t:Q_t(X_0)\to X_0$ satisfying
\[
D_t\circ Q_t=\operatorname{id}_{X_0}.
\]
It recovers the actual upper input.
\end{definition}

\begin{definition}[Autonomous reduced evolution]
An autonomous reduced evolution is a map $F_t:Y_0\to Y_t$ satisfying
\[
F_t\circ P_0=P_t\circ\Phi_t.
\]
\end{definition}

\begin{definition}[Effective merger]
If $F_t$ exists, an effective merger occurs when distinct $y,y'\in Y_0$
satisfy $F_t(y)=F_t(y')$.
\end{definition}

\begin{theorem}[Projection--descent and recovery]\label{thm:descent-recovery}
For the typed data above:
\begin{enumerate}
\item a set-theoretic representative section exists when $Q_t$ is surjective
and the relevant choice principle is available; measurable, continuous,
local, or Lipschitz sections require corresponding selection theorems;
\item an exact decoder exists if and only if $Q_t$ is injective, in which
case it is the inverse of $Q_t$ on its image;
\item an autonomous reduced evolution exists if and only if
\begin{equation}
P_0(x)=P_0(x')
\quad\Longrightarrow\quad
P_t(\Phi_t x)=P_t(\Phi_t x')
\label{eq:descent}
\end{equation}
for all $x,x'\in X_0$, and it is then unique;
\item if~\eqref{eq:descent} holds and two distinct initial effective states
have the same final image, then $F_t$ is noninjective and the prior effective
state cannot be decoded uniquely from the final one.
\end{enumerate}
\end{theorem}

\begin{proof}
The first item is the definition of a section of a surjection.  If
$D_tQ_t=\operatorname{id}$ and $Q_t(x)=Q_t(x')$, applying $D_t$ gives
$x=x'$.  Conversely, an injective map has an inverse on its image.  For the
third item, necessity follows by applying $F_t$ to equal $P_0$-images.
Under~\eqref{eq:descent}, define
\[
F_t(P_0x):=P_t(\Phi_tx).
\]
The implication makes this independent of the representative, and
surjectivity of $P_0$ onto $Y_0$ gives existence and uniqueness.  The final
item is immediate from the definition of injectivity.
\end{proof}

\begin{corollary}[Noninjectivity is not a right-section obstruction]
Noninjectivity of $Q_t$ rules out exact microscopic recovery, not
representative selection.  Surjectivity, together with the required
regularity category, is the gate for a right section.
\end{corollary}

\subsection{Admissible domains and explicit margins}

\begin{definition}[Declared margins]
Let $m_1,\ldots,m_r$ be continuous real-valued margins on a metric upper
domain.  They may encode such requirements as operator-domain control,
spectral separation, bounded leakage, truncation error, section conditioning,
or physical hyperbolicity.  For $\delta\ge0$, set
\[
A^\delta:=\{x:m_j(x)\ge\delta\text{ for every }j\}.
\]
\end{definition}

\begin{definition}[Admissible domain]
A subset $A\subset X_0$ is admissible for a declared protocol family
$\Pi_A$ if:
\begin{enumerate}
\item $A\subset A^\delta$ for some $\delta>0$ and $P_0|_A$ is measurable;
\item for each $\pi\in\Pi_A$ there is a representative section
$S_{A,\pi}:P_0(A)\to A$ of $P_0|_A$ in the declared regularity category;
\item a declared upper step $R_{A,\pi}:A\to X_0$ induces a well-typed
section-dependent self-map
\[
T_{A,\pi}:=P_0\circ R_{A,\pi}\circ S_{A,\pi}
   :P_0(A)\longrightarrow P_0(A);
\]
\item the fixed-point, encoding, or approximation contract invoked on $A$
is stated explicitly and verified on its own invariant domain.
\end{enumerate}
\end{definition}

\begin{axiom}[Finite admissibility]
For the declared theory class, there is no single domain $A=X_0$ and
protocol family for which all required positive-margin, regularity, and
closure contracts hold.  Equivalently, no one globally valid admissible
encoding exists at this adopted axiom tier.
\end{axiom}

\begin{remark}
Finite admissibility excludes a single globally valid \emph{admissible
encoding} by definition.  It does not, by itself, exclude a global relation,
a set-theoretic section, or even an autonomous reduced map.
\end{remark}

\subsection{Basin-local Fundamental Contractivity Condition}

\begin{definition}[Exact basin-local FCC]\label{def:fcc}
Fix $A$, $\pi$, and $T:=T_{A,\pi}$.  A basin contract is a nonempty complete
metric subspace $D_\alpha\subset P_0(A)$ such that
\[
T(D_\alpha)\subseteq D_\alpha,
\qquad
d_Y(Tu,Tv)\le q_\alpha d_Y(u,v)
\]
for all $u,v\in D_\alpha$, where $0\le q_\alpha<1$.  FCC means this exact,
basin-local contract, not one contraction on a union of outcome basins.
\end{definition}

\begin{theorem}[Basin-local fixed point]\label{thm:basin-fixed-point}
Every basin $D_\alpha$ satisfying Definition~\ref{def:fcc} contains a unique
fixed point $y_\alpha^\ast$, and
\[
d_Y(T^ny,y_\alpha^\ast)\le
q_\alpha^n d_Y(y,y_\alpha^\ast)
\quad(y\in D_\alpha).
\]
\end{theorem}

\begin{proof}
This is Banach's fixed-point theorem on the complete invariant domain
$D_\alpha$.
\end{proof}

\begin{definition}[Approximate orbit contract]
On an invariant domain $D$, an approximate orbit contract is
\[
d_Y(Tu,Tv)\le\kappa d_Y(u,v)+c\varepsilon,
\qquad 0\le\kappa<1.
\]
\end{definition}

\begin{proposition}[Approximate pairwise estimate]\label{prop:approximate}
If two forward orbits remain in $D$, then
\[
d_Y(T^nu,T^nv)
\le \kappa^n d_Y(u,v)
 \frac{1-\kappa^n}{1-\kappa}\,c\varepsilon.
\]
This estimate alone neither makes $T$ a Banach contraction nor proves that
$T$ has a fixed point.
\end{proposition}

\begin{proof}
Iterate the affine recurrence for the distance.
\end{proof}

\subsection{Fixed points and basins}

\begin{definition}[Basin]
A basin is an invariant domain $D_\alpha$ equipped with either the exact FCC
contract of Definition~\ref{def:fcc}, or a separately stated stability
theorem.  Distinct outcome basins are not combined into one Banach domain.
\end{definition}

\subsection{Coherence capacity}

\begin{definition}[Coherence capacity]
For declared margins $m_j$, define the diagnostic
\[
C(x):=\min_j m_j(x).
\]
Thus $C(x)>0$ records positive slack in every declared condition and
$C(x)=0$ records saturation of at least one condition.
\end{definition}

\subsection{Admissibility barriers}

\begin{definition}[Margin boundary and barrier]\label{def:admissibility-barrier}
For a declared contract, let
\[
X_{\rm adm}:=\{x:C(x)>0\},
\qquad
\partial X_{\rm adm}:=\overline{X_{\rm adm}}
   \setminus\operatorname{int}(X_{\rm adm}).
\]
A pathwise admissibility barrier is a subset
$\mathcal B\subseteq\partial X_{\rm adm}$ across which a specified positive
margin cannot be continued.  The failed margin must be named.
\end{definition}

\begin{definition}[Section-conditioning diagnostic]
For Lipschitz representative sections of $Q_t$, define
\[
\kappa_{\rm sec}(t):=
\inf\{\operatorname{Lip}(S_t):Q_tS_t=\operatorname{id}_{Y_t}\},
\]
with $\kappa_{\rm sec}(t)=+\infty$ if no Lipschitz section exists.
\end{definition}

\begin{remark}
Blow-up of $\kappa_{\rm sec}$, loss of surjectivity, or failure of bounded
Riesz-projector continuation can supply explicit barrier margins in a
realization.  None is assumed merely to manufacture an obstruction, and none
alone is a physical singularity or state-selection law.
\end{remark}

\subsection{Invariant measure (conditional)}

\begin{assumption}[Measure data, only when invoked]
For probabilistic statements, supply either a preparation probability
measure $\mu$ on the relevant upper domain or a stationary probability
measure $\nu$ for a declared self-map.  Invariance means
$R_\#\mu=\mu$ for $R:A\to A$, or $T_\#\nu=\nu$ for
$T:Y_A\to Y_A$; it is not meaningful without a same-space map.
\end{assumption}

\subsection{Encodings}

\begin{definition}[Encoding]
An encoding is a tuple
\[
\mathcal E=(A,\;Z,\;E,\;F,\;\varepsilon)
\]
where $A$ is admissible, $Z$ is a normed space,
$E:P_0(A)\to Z$, and
\[
E(T_{A,\pi}(y)) = F(E(y)) + \delta(y),
\quad \|\delta(y)\|\le c\,\varepsilon.
\]
\end{definition}
\begin{remark}
Coherence capacity $C(x)$ is not an observable, not a conserved quantity,
and not a dynamical degree of freedom. It is a diagnostic margin (a stability/condition
number) organizing when admissible encodings exist and when they fail.
\end{remark}

\subsection{Interpretation postulate}

\begin{postulate}
When an admissible encoding exists, its reduced structure
may be interpreted as a physical description.
\end{postulate}
\begin{remark}
Local sections $S_{A,\pi}$ may exist on admissible domains.
Whether they extend is a separate selection and conditioning problem.
Thus local representability does not imply stable global representability,
but neither does it logically forbid a set-theoretic global section.
\end{remark}

\subsection{Failure modes}

The corresponding conclusions are unavailable if:
\begin{itemize}
\item no complete invariant basin satisfies FCC where a fixed point is claimed;
\item no invariant measure exists where probability is claimed;
\item no representative section exists in the required regularity category;
\item the declared margins, factorization, or robustness estimates fail.
\end{itemize}


\section{Core Structural Results}

Fix an admissible domain $A$, protocol $\pi$, and the well-typed self-map
$T:=T_{A,\pi}$.

\subsection{Basin-local stability}

\begin{theorem}[Universality within one exact basin]
If $D_\alpha$ satisfies exact basin-local FCC, then every two orbits in
$D_\alpha$ converge to the same unique fixed point:
\[
\lim_{n\to\infty}d_Y(T^ny,T^ny')=0
\qquad(y,y'\in D_\alpha).
\]
\end{theorem}

\begin{proof}
Apply Theorem~\ref{thm:basin-fixed-point} to both orbits, or use
$d_Y(T^ny,T^ny')\le q_\alpha^n d_Y(y,y')$.
\end{proof}

\begin{remark}[No cross-basin conclusion]
Separate invariant basins may contain separate fixed points.  Neither exact
FCC nor Proposition~\ref{prop:approximate} may be applied to their union
unless that union independently satisfies the required complete invariant
contract.
\end{remark}
\subsection{Admissible prediction depth}

\begin{definition}[Admissible prediction depth]
Fix an admissible domain $A$ and protocol $\pi\in\Pi_A$.
For $x\in A$, define the admissible prediction depth
\[
\begin{aligned}
N_{\max}(x;A,\pi,\varepsilon)
:=\sup\big\{n\in\mathbb{N}\ \big|\ &
T_{A,\pi}^k(P_0(x))\in P_0(A)\\
&\text{for every }0\le k\le n\big\}.
\end{aligned}
\]
\end{definition}

\begin{remark}
$N_{\max}$ is an iteration-count diagnostic.  Its relation to physical time
or to the size of a margin must be proved in a realization.
\end{remark}
\subsection{Computational limits (optional)}

\begin{remark}
Finite, instance-dependent prediction depth does not imply undecidability.
An undecidability result requires a specified input language and a reduction
from a known undecidable problem to a declared reachability or admissibility
question.  A0 supplies no such reduction.
\end{remark}

\subsection{Forgetting of incoherent distinctions}

\begin{corollary}[Basin-local metric forgetting]
For $y,y'$ in one exact FCC basin,
\[
d_Y(T^ny,T^ny')\longrightarrow0.
\]
Under only the approximate orbit contract,
\[
\limsup_{n\to\infty}d_Y(T^ny,T^ny')
\le \frac{c\varepsilon}{1-\kappa},
\]
provided both orbits remain in the same invariant domain.  This is a metric
statement; information or entropy claims require additional structures.
\end{corollary}

\subsection{A valid reduced-self-map obstruction}

\begin{theorem}[Finite-diameter contraction obstruction]
\label{thm:diameter-obstruction}
Let $(Y_A,d)$ have finite positive diameter $D$, and let
$G:Y_A\to Y_A$ satisfy
\[
d(Gy,Gy')\le\kappa d(y,y')+c\varepsilon,
\qquad0\le\kappa<1.
\]
Then
\[
\operatorname{diam}G(Y_A)\le\kappa D+c\varepsilon.
\]
If $(1-\kappa)D>c\varepsilon$, then $G$ is not surjective and has no
right section $S:Y_A\to Y_A$ with
$G\circ S=\operatorname{id}_{Y_A}$.
\end{theorem}

\begin{proof}
Take the supremum of the displayed estimate over all pairs.  Under the
strict inequality, the image diameter is smaller than $D$, whereas a
surjective image would be all of $Y_A$ and would have diameter $D$.
\end{proof}

\begin{remark}
Theorem~\ref{thm:diameter-obstruction} concerns a self-map of one reduced
space.  It cannot be applied to a cross-level map merely because the latter
is noninjective.
\end{remark}

\subsection{Recovery and irreversibility are separate}

\begin{proposition}[Recovery classification]
For the typed maps of Theorem~\ref{thm:descent-recovery}:
\begin{enumerate}
\item noninjectivity of $Q_t$ obstructs exact upper recovery;
\item noninjectivity of an existing $F_t$ obstructs recovery of the prior
effective state;
\item neither fact alone proves a physical arrow of time.
\end{enumerate}
A physical arrow additionally needs oriented dynamics and an asymmetric
property such as a noninvertible physical semigroup, a monotone Lyapunov or
entropy functional, or a boundary condition.
\end{proposition}

\subsection{Measure-dependent reduced probability}

\begin{theorem}[Measure-dependent reduced kernel]\label{thm:reduced-kernel}
Let $A$ and $P_0(A)$ be standard Borel spaces, let $\mu$ be a probability
measure on $A$, and let $\{\mu_y\}$ be a regular conditional distribution of
$x$ given $P_0(x)=y$.  For a measurable upper step $R:A\to X_0$, define
\[
K(y,B):=\mu_y\{x:P_0(Rx)\in B\}.
\]
Then $K$ is a Markov kernel, up to the usual $(P_0)_\#\mu$-null sets.  If
$P_0R$ descends to $F$ on $P_0(A)$, then
$K(y,\cdot)=\delta_{F(y)}$ almost everywhere.
\end{theorem}

\begin{proof}
Regular conditional distributions exist on standard Borel spaces.
Measurability and countable additivity pass through the measurable preimage.
Under autonomous descent, $P_0(Rx)$ is constant on each conditional fiber.
\end{proof}

\begin{definition}[Basin weights]
Let $\{B_i\}$ be a measurable partition of $P_0(A)$ and let $\mu$ be a
declared probability measure on $A$.  Define
\[
W_i :=
\mu\big(P_0^{-1}(B_i)\cap A\big),
\]
or condition and renormalize this expression on a declared selection event.
\end{definition}

\begin{remark}
The weights are probabilities because $\mu$ was supplied, not because the
projection has fibers.  Different upper measures can induce different
weights and kernels.  No Born rule is asserted.
\end{remark}


\subsection{Re-encoding criterion on overlaps}

\begin{theorem}[Exact encoding factorization]\label{thm:encoding-factor}
Let $E_1:Y\to Z_1$ and $E_2:Y\to Z_2$.  A map
$f:E_1(Y)\to Z_2$ satisfying $E_2=f\circ E_1$ exists if and only if
\[
E_1(y)=E_1(y')\quad\Longrightarrow\quad E_2(y)=E_2(y')
\]
for all $y,y'\in Y$.  When it exists, $f$ is unique on $E_1(Y)$.
\end{theorem}

\begin{proof}
Necessity is immediate.  For sufficiency define
$f(E_1(y)):=E_2(y)$; the implication makes this well-defined.
\end{proof}

\begin{remark}
Overlap of domains alone does not produce a re-encoding.  Approximate
factorization requires a quantitative fiber-consistency estimate and a
controlled extension in the declared regularity class.  A0 therefore makes
no universal maximal-content claim.
\end{remark}

\subsection{Finite admissibility and global scope}

\begin{proposition}[No global admissible encoding at the adopted axiom tier]
Under the finite-admissibility axiom, no object satisfying the definition of
an admissible encoding has domain $X_0$.
\end{proposition}

\begin{proof}
This is the direct content of the axiom and the definition of an encoding.
\end{proof}

\begin{remark}
This proposition does not prohibit a global relation or a global reduced map
that fails one of the adopted admissibility margins.
\end{remark}

% =====================================================
\section*{Part II: Encodings and Boundaries}
\addcontentsline{toc}{section}{Part II: Encodings and Boundaries}
\section{Encoding Atlas}

This section lists admissible encodings as local charts on neighborhoods
of stable basins.  Fix $A$ and $\pi$ and write
\[
Y_A:=P_0(A),
\qquad T_A:=T_{A,\pi}:Y_A\to Y_A .
\]
Each encoding is specified by sufficient inequality contracts involving
the contraction constant $\kappa$, tolerance $\varepsilon$, and
coherence capacity $C$.
\begin{remark}[Contract nature of atlas inequalities]
The inequalities in this atlas are \emph{sufficient contracts} guaranteeing that an encoding closes
to tolerance on the stated domain. They are not classification theorems: failure of a listed
inequality does not imply that no encoding exists.
\end{remark}
\begin{remark}[Status of numerical thresholds]
Numerical coefficients displayed below are declared chart-design thresholds,
not universal constants derived by A0.  A physical realization must justify
or replace them and verify every associated closure estimate.
\end{remark}

\subsection{Chart-local control quantities}

The following quantities are \emph{chart-local}:
they are defined only when the corresponding encoding supplies
the structures needed to define them.
They are not global invariants of the typed projection system.

\begin{itemize}
\item Basin separation scale (when a basin partition exists):
\[
\Delta_B := \inf_{i\neq j} d_Y(B_i,B_j).
\]

\item Capacity floor on the chart domain:
\[
C_{\min} := \inf_{x\in A} C(x).
\]

\item Contractivity slack:
\[
s_\kappa := 1-\kappa.
\]
\end{itemize}

All inequalities below are sufficient but not necessary.

%-------------------------------------------------

\subsection{Encoding \texorpdfstring{E$_1$}{E1}: Finite-Basin Operator Chart}

\paragraph{Existence hypothesis}
The reduced space $Y_A$ admits a finite or effectively finite
collection of basins $\{B_i\}$ with well-defined separation.

\paragraph{Validity conditions}
\begin{align}
\varepsilon &\le \tfrac{1}{10}\Delta_B, \\
\kappa &\le \tfrac{1}{2}, \\
C_{\min} &\ge C_1 > 0.
\end{align}

\paragraph{Closure contract}
The chart is admitted only if a linear map $F$ on $Z$ is supplied and
verified to satisfy
\[
E(T_A(y)) = F(E(y)) + \delta(y),
\quad \|\delta(y)\|\le c_1\,\varepsilon .
\]

\paragraph{Failure}
This chart contract fails when basin separation is lost, $\kappa\to 1$,
or $C_{\min}\to 0$; another encoding may still exist.

%-------------------------------------------------

\subsection{Encoding \texorpdfstring{E$_2$}{E2}: Local Patch Algebra Chart}

\paragraph{Existence hypothesis}
The admissible domain admits a cover by overlapping admissible
subdomains $\{A_\alpha\}$, but no global admissible section exists.

\paragraph{Validity conditions}
\begin{align}
\inf_\alpha \inf_{x\in A_\alpha} C(x) &\ge C_2 > 0, \\
\sup_\alpha \kappa_\alpha &\le \kappa_2 < 1.
\end{align}

On overlaps $A_{\alpha\beta}$,
\[
\|E_\beta(y) - f_{\alpha\beta}(E_\alpha(y))\|
\le c_2\,\varepsilon .
\]

\paragraph{Obstruction contract}
A claimed global obstruction must be exhibited by incompatible transition
cocycles or a proved lower bound forcing every glue error above
$\varepsilon$.  Local overlap alone is not an obstruction.

\paragraph{Failure}
Overlap incompatibility or capacity collapse.

%-------------------------------------------------

\subsection{Encoding \texorpdfstring{E$_3$}{E3}: Two-Derivative Bookkeeping Chart}

\paragraph{Existence hypothesis}
Variations of basin representatives are slow enough that
a second-order truncation is meaningful.

\paragraph{Validity conditions}
\begin{align}
C_{\min} &\ge C_3 > 0, \\
q_\alpha &\le 1-s_3<1
\end{align}

Let $\eta_3$ denote the ratio of neglected higher-order terms
to retained second-order terms.
Require
\[
\eta_3 \le c_3\,\varepsilon .
\]

\paragraph{Closure contract}
One must verify
\[
E(T_A(y)) = E(y) + \mathcal{D}[E(y)] + \delta(y),
\quad \|\delta(y)\|\le c_3'\,\varepsilon .
\]

\paragraph{Failure}
Higher-order terms become unsuppressed or $C\to 0$.

%-------------------------------------------------

\subsection{Encoding \texorpdfstring{E$_4$}{E4}: Scale-Step Truncation Chart}

\paragraph{Existence hypothesis}
There exists a hierarchy of reduced coordinates ordered by relevance.

\paragraph{Validity conditions}
\begin{align}
C_{\min} &\ge C_4 > 0, \\
\rho_N &\le \varepsilon, \\
\kappa &\le 1 - c_4\,\varepsilon .
\end{align}

Here $\rho_N$ denotes a chart-local bound on the cumulative
contribution of discarded coordinates.

\paragraph{Closure contract}
One must verify
\[
E(T_A(y)) = F(E(y)) + \delta(y),
\quad \|\delta(y)\|\le c_4'\,\varepsilon .
\]

\paragraph{Failure}
Discarded contributions exceed tolerance or $\kappa\to 1$.

%-------------------------------------------------

\subsection{Encoding \texorpdfstring{E$_5$}{E5}: Auxiliary Fixed-Point Chart}

\paragraph{Existence hypothesis}
There exists an auxiliary representation $Z_{\mathrm{aux}}$
in which the reduced dynamics is near a stable fixed point.

\paragraph{Validity conditions}
\begin{align}
\|\beta(E(y))\| &\le c_5\,\varepsilon, \\
\|D\beta\| &\le 1 - s_5, \\
C_{\min} &\ge C_5 > 0.
\end{align}

\paragraph{Closure contract}
The auxiliary map must preserve a complete neighborhood and satisfy an exact
contraction there, or a separate fixed-point existence theorem must be
provided.  A spectral-radius bound alone is insufficient for a nonnormal
linearization.

\paragraph{Failure}
Loss of auxiliary stability or capacity collapse.

%-------------------------------------------------

\subsection{Encoding \texorpdfstring{E$_6$}{E6}: Spectral Operator Chart}

\paragraph{Existence hypothesis}
The encoding supplies operators whose spectra are meaningful
on the admissible domain.

\paragraph{Validity conditions}
\begin{align}
\varepsilon &\le \tfrac{1}{10}\Delta_D, \\
\|[D,a]\| &\le K_6, \\
C_{\min} &\ge C_6 > 0.
\end{align}

\paragraph{Closure contract}
Dependence on spectral data only, with its error bound, must be verified for
the declared operator domain.

\paragraph{Failure}
Spectral gap closure or operator norm divergence.

%-------------------------------------------------

\subsection{Encoding \texorpdfstring{E$_7$}{E7}: Boundary and Reset Chart}

\paragraph{Existence hypothesis}
Capacity reaches a declared boundary and an additional continuation or reset
rule $J_e$ is supplied.

\paragraph{Validity conditions}
\begin{align}
0 < C_{\min} &\le C_7, \\
\tau_e &\ge \tau_7 > 0.
\end{align}

Reset labels must be stable under $\varepsilon$-perturbations, $J_e$ must be
measurable, and any required conservation law must be checked.  Post-reset
basins must satisfy FCC with positive slack.  If probabilities are assigned,
a measure or stochastic kernel must also be supplied.

\paragraph{Failure}
Event clustering, unstable post-selection basins,
or $C<0$.
%-------------------------------------------------
\subsection{Encoding \texorpdfstring{E$_8$}{E8}: Relational Transition Chart}

This encoding applies when no single-valued reduced evolution
is admissible, but stable relational structure persists.

\paragraph{Existence hypothesis}
\begin{itemize}
\item The descent criterion~\eqref{eq:descent} fails on the neighborhood of
interest, so no autonomous single-valued reduced map exists there.
\item The typed reduced relation
\[
\mathcal{R}_t :=
\{(P_0x,P_t\Phi_tx):x\in X_0\}\subseteq Y_0\times Y_t
\]
admits a stable restriction.
\end{itemize}

\paragraph{Variables}
The reduced data are not points in a coordinate space, but
\emph{relations}:
\[
Z_{\mathrm{rel}} \subseteq Y_0\times Y_t.
\]

\paragraph{Chart map}
\[
E : Y_0 \to \mathcal{P}(Y_t),
\qquad
E(y) := \{y' \mid (y,y')\in\mathcal{R}_t\},
\]
where $\mathcal{P}(Y_t)$ denotes the power set of $Y_t$.

\paragraph{Closure contract}
For a declared semigroup $\Phi_{s+t}=\Phi_s\Phi_t$ with compatible
time-indexed reductions, relational composition must be checked in the typed
form
\[
\mathcal{R}_s\circ\mathcal{R}_t
\subseteq \mathcal{R}_{s+t}
\quad\text{up to the declared tolerance}.
\]
Transitivity of one fixed-time relation is not automatic.

\paragraph{Validity conditions}
\begin{align}
C_{\min} &\ge C_8 > 0, \\
\text{Autonomous descent} &\text{ fails on the domain.}
\end{align}

\paragraph{Failure}
\begin{itemize}
\item Relational structure becomes unstable or dense;
\item No controlled restriction of $\mathcal{R}_t$ exists;
\item Coherence capacity collapses globally.
\end{itemize}

\paragraph{Remarks}
This encoding is one possible weak description when deterministic descent
fails.  It captures persistent constraints on
allowed transitions without introducing local coordinates, states,
or probabilities.
%-------------------------------------------------
\subsection{Encoding \texorpdfstring{E$_9$}{E9}: Statistical Ensemble Chart}

This encoding applies when local or pointwise reduced descriptions fail,
but stable statistical structure persists without a well-defined
invariant measure.

\paragraph{Existence hypothesis}
\begin{itemize}
\item No admissible invariant or stationary measure exists on the
      admissible domain.
\item Long-time empirical averages or ensemble summaries stabilize
      under reduced evolution.
\item Single-trajectory descriptions do not close, but aggregate
      quantities do.
\end{itemize}

\paragraph{Variables}
The reduced variables are ensemble summaries:
\[
Z_{\mathrm{stat}} := \{\text{empirical distributions, moments, or coarse summaries}\}.
\]

\paragraph{Chart map}
\[
E : Y_A \to Z_{\mathrm{stat}},
\qquad
E(y) := \lim_{N\to\infty} \frac{1}{N}\sum_{k=1}^N \mathcal{O}(T_A^k(y)),
\]
whenever the limit exists for the chosen summary observable $\mathcal{O}$.

\paragraph{Closure contract}
The chart is admitted only if ensemble evolution closes approximately:
\[
E(T_A(y)) = F(E(y)) + \delta(y),
\qquad \|\delta(y)\|\le c_9\,\varepsilon,
\]
even though no pointwise closure exists.

\paragraph{Validity conditions}
\begin{align}
C_{\min} &\ge C_9 > 0, \\
\text{No invariant measure } \mu &\text{ exists on the domain.}
\end{align}

\paragraph{Failure}
\begin{itemize}
\item Ensemble summaries fail to stabilize;
\item Strong nonstationarity destroys aggregate closure;
\item Capacity collapses and no admissible encoding remains.
\end{itemize}

\paragraph{Remarks}
This empirical-summary encoding does not itself invoke probability, random
variables, or stochastic dynamics.  It is not universally ordered above or
below the relational encoding; comparison requires a proved re-encoding.

%-------------------------------------------------

\subsection{Atlas summary}

Admissible encodings range from coordinate descriptions to relational and
statistical summaries.  In a particular realization one may observe a
schematic transition
\[
E_1 \;\to\; \cdots \;\to\; E_7 \;\to\; E_9 \;\to\; E_8,
\]
but A0 does not prove this ordering.  Each arrow requires an overlap
factorization and a verified margin comparison.


\section{Admissibility Boundaries and Selection}

This section analyzes the failure of admissible reduced descriptions.
All notions refer to the definitions in Section~1 and are purely structural.

\subsection{Admissibility barriers}

We recall the definition of an admissibility barrier from
Definition~\ref{def:admissibility-barrier}.
No additional structure is assumed here.

\subsection{Termination of encodings}

\begin{proposition}[Loss of an encoding certificate at a barrier]
Let $\mathcal{E}=(A,Z,E,F,\varepsilon)$ be an admissible encoding.
If a trajectory reaches a barrier at which one of $\mathcal E$'s required
margins vanishes, then the existing certificate for $\mathcal E$ ends there.
Continuation requires a new proof restoring that margin, a different
encoding, or an additional continuation law.
\end{proposition}

\begin{proof}
Positive slack in every declared margin is part of admissibility.  At the
named zero margin, that hypothesis is no longer available.
\end{proof}

\subsection{Selection events}

\begin{definition}[Declared selection or reset event]
A selection event consists of a margin crossing together with a supplied
continuation rule
\[
J_e:\mathcal B_e\times\Lambda_e\longrightarrow
\bigcup_\alpha D_\alpha,
\]
where $\Lambda_e$ contains any extra label or random input.  The rule must be
typed, measurable in probabilistic uses, and compatible with every claimed
conservation law.
\end{definition}

Margin saturation alone terminates a certificate; it does not choose an
outcome or construct $J_e$.

\subsection{Irreversibility of selection}

\begin{proposition}[Conditional effective irreversibility]
A declared reset is irreversible at the effective level if its induced map
from pre-event effective states to post-event effective states is
noninjective.  Exact upper recovery is impossible if the corresponding typed
upper-to-output map is noninjective.  Neither statement follows from loss of
FCC alone, and neither establishes a physical arrow of time without an
oriented asymmetric evolution law.
\end{proposition}

\subsection{Boundary layers}

\begin{definition}[Boundary layer]
A $\delta$-boundary layer for the declared margins is
\[
\mathcal L_\delta:=\{x:0<C(x)\le\delta\}.
\]
\end{definition}

\begin{remark}
Small margin identifies proximity to failure only for the declared contract.
Sensitivity amplification, coordinate independence, or common scaling
between encodings requires a comparison theorem for their margins.
\end{remark}

\subsection{Records}

\begin{definition}[Record]
A record for $T$ is a measurable map $R_{\rm rec}$ on a forward-invariant
post-event domain such that
\[
R_{\rm rec}\circ T=R_{\rm rec}.
\]
\end{definition}

\begin{proposition}[Record persistence]
Once such an $R_{\rm rec}$ is supplied, its value is constant along every
forward orbit in its domain.
\end{proposition}

\begin{proof}
Iterate the defining identity.
\end{proof}

FCC alone does not construct a nontrivial record.  A partial order on events
also requires an oriented event relation and record compatibility.

\subsection{Probability at boundaries (conditional)}

Let $\mu$ be a declared upper probability measure and let
$\{B_i\}$ be a measurable partition of the post-event reduced domain.  The
conditional numbers
\[
W_i =
\frac{\mu(P_0^{-1}(B_i)\cap A_e)}
{\sum_j \mu(P_0^{-1}(B_j)\cap A_e)}
\]
are probabilities when the denominator is positive.  They describe the
chosen measure and event domain $A_e$; they are not selected by the barrier
and are not asserted to be Born weights.

\subsection{No global reversibility}

\begin{remark}
Finite admissibility rules out one globally certified admissible encoding at
the adopted axiom tier.  It does not by itself rule out an invertible global
relation, a reversible upper evolution, or a different reduced description
whose hypotheses are not those of A0.
\end{remark}

\subsection{Boundary summary}

Admissibility barriers mark the limits of a named certificate.  At such a
limit one must stop, prove a continuation, switch to another verified
encoding, or supply a reset law.  Selection, probability, records, and
irreversibility are separate gates.
% =====================================================
\section*{Part III: Structural Closure}
\addcontentsline{toc}{section}{Part III: Structural Closure}
\section{Exact and Controlled Relations Between Encodings}

Category language is exact.  We therefore define an exact category first
and keep tolerance errors as explicit controlled data rather than silently
quotienting by a relation that may fail to be a congruence.

\subsection{The exact local category}

Fix one admissible domain $A$, protocol $\pi$, reduced space
$Y_A=P_0(A)$, and self-map $T_A$.

\begin{definition}[Exact encoding object]
An object of $\mathbf{AdmEnc}_0(A,T_A)$ is a triple
\[
\mathcal E=(Z,E,F)
\]
with $E:Y_A\to Z$, $F:Z\to Z$, and
\[
E\circ T_A=F\circ E.
\]
\end{definition}

\begin{definition}[Exact re-encoding morphism]
A morphism $f:\mathcal E_1\to\mathcal E_2$ is a map
$f:Z_1\to Z_2$ satisfying
\[
E_2=f\circ E_1,
\qquad
f\circ F_1=F_2\circ f
\quad\text{on }E_1(Y_A).
\]
\end{definition}

\begin{proposition}
Identity maps and ordinary composition make
$\mathbf{AdmEnc}_0(A,T_A)$ a category.
\end{proposition}

\begin{proof}
Both factorization identities are preserved by identity maps and by
composition.
\end{proof}

\subsection{The correctly oriented universal object}

\begin{definition}[Identity encoding]
\[
\mathcal C_A:=(Y_A,\operatorname{id}_{Y_A},T_A).
\]
\end{definition}

\begin{theorem}[Local initiality]\label{thm:initiality}
$\mathcal C_A$ is initial in $\mathbf{AdmEnc}_0(A,T_A)$.  For every exact
encoding $\mathcal E=(Z,E,F)$, the unique morphism
$\mathcal C_A\to\mathcal E$ is $E$.
\end{theorem}

\begin{proof}
The object identity $E\circ T_A=F\circ E$ is exactly the dynamical
intertwining condition.  The equation $E=f\circ\operatorname{id}$ forces
$f=E$.
\end{proof}

\begin{proposition}[Recovery from an encoding]
A morphism $\mathcal E\to\mathcal C_A$ exists if and only if
$\operatorname{id}_{Y_A}$ factors through $E$.  Equivalently,
\[
E(y)=E(y')\Longrightarrow y=y',
\]
and a suitable inverse on $E(Y_A)$ exists in the declared regularity
category.
\end{proposition}

\begin{remark}
Version~1 reversed this arrow and called $\mathcal C_A$ terminal.  A general
encoding can discard reduced information, so recovery of $Y_A$ from its
coordinates is an additional injectivity and regularity theorem.
\end{remark}

\subsection{Controlled approximate re-encodings}

For approximate encodings, retain an explicit defect
\[
\|E\circ T_A-F\circ E\|\le\eta
\]
on a declared domain.  A controlled arrow $f:\mathcal E_1\to\mathcal E_2$
records both a factorization defect
$\|E_2-fE_1\|\le\eta_f$ and a dynamical-intertwining defect.  If $g$ is
$L_g$-Lipschitz, then the factorization defect of $g\circ f$ is at most
\[
\eta_g+L_g\eta_f.
\]
Thus controlled arrows compose with an explicit error budget.  A quotient
category may be formed only after the proposed tolerance relation is proved
to be an equivalence relation compatible with this composition; A0 does not
assume such a quotient.

\subsection{Overlaps and global scope}

For encodings on $A_1$ and $A_2$, first restrict to
$Y_{12}:=P_0(A_1\cap A_2)$.  Theorem~\ref{thm:encoding-factor}, plus the
declared regularity and error estimates, decides whether a transition map
exists.  Nonempty overlap alone is insufficient.

The finite-admissibility axiom excludes an exact or controlled
\emph{admissible} object whose domain is all of $X_0$.  A particular pair of
charts fails to glue only when a transition obstruction or incompatible
margin is actually proved.

\subsection{Pushforward probability as an optional functor}

Let $\nu$ be a declared probability measure on $Y_A$, and restrict to
measurable exact encodings and measurable morphisms.  Assign to
$\mathcal E=(Z,E,F)$ the pushforward probability space
$(Z,E_\#\nu)$, and to $f:\mathcal E_1\to\mathcal E_2$ the measurable
pushforward map $f$.  Since $E_2=fE_1$,
\[
(E_2)_\#\nu=f_\#(E_1)_\#\nu,
\]
so this assignment is a functor to probability spaces.  The functor uses
the supplied $\nu$; category structure does not create it.

\subsection{Category summary}

The identity encoding is locally initial, not terminal.  Arrows out of it
are encodings; arrows back to it are recovery maps.  Approximate arrows carry
error budgets, and global gluing remains a separate descent problem.

\section*{Part IV: Total Atlas}
\addcontentsline{toc}{section}{Part IV: Total Atlas}

\section{Total Atlas: A Global Chart-of-Charts}

This section introduces a global object that \emph{contains} all admissible encodings
without asserting the existence of any single globally valid encoding.
It formalizes the idea of a \emph{total chart} as a chart-of-charts.

\subsection{Local encoding fiber}

\begin{definition}[Encoding fiber]
For each $x\in X_0$, define
\[
\mathrm{Enc}(x):=
\{\mathcal E:\mathcal E\text{ is a verified exact or controlled encoding
on some }A\ni x\}.
\]
\end{definition}

\begin{remark}
$\mathrm{Enc}(x)$ is the set of all reduced descriptions available at $x$.
Finite admissibility says no single verified chart has domain $X_0$; it does
not imply that any point is uncovered by the union of local charts.
\end{remark}

\subsection{Total encoding family}

\begin{definition}[Total encoding projection]
Define the total encoding space
\[
\mathbb{E} := \{(x,\mathcal{E}) \mid x\in X_0,\ \mathcal{E}\in \mathrm{Enc}(x)\},
\]
with projection map
\[
\pi:\mathbb{E}\to X_0,\qquad \pi(x,\mathcal{E})=x.
\]
\end{definition}

\begin{remark}
The fiber $\pi^{-1}(x)$ is canonically identified with $\mathrm{Enc}(x)$.
Without a topology and local trivializations, $\pi$ is a set-theoretic
fibered family, not a fiber bundle.  It is not itself an encoding.
\end{remark}

\subsection{Admissibility limits as fiber collapse}

\begin{definition}[Encodability set]
Define the encodable subset of $X_0$ by
\[
X_{\mathrm{enc}} := \{x\in X_0 \mid \mathrm{Enc}(x)\neq\varnothing\}.
\]
\end{definition}

\begin{definition}[Encoding boundary set]
Define the topological boundary as
\[
\partial X_{\mathrm{enc}} :=
\overline{X_{\mathrm{enc}}}
\setminus\operatorname{int}(X_{\mathrm{enc}}).
\]
\end{definition}

\begin{remark}
Boundary points are accumulation points of both encodable and non-interior
behavior.  A claim that availability actually collapses there requires
upper-semicontinuity or a specific margin theorem.
\end{remark}

\subsection{Chart viability field}

The atlas in Section~3 provides \emph{sufficient contracts} in terms of local control
quantities. We now compress those contracts into a single diagnostic.

\begin{definition}[Chart viability functional]
For an encoding $\mathcal{E}$ and a point $x\in X_0$, define a chart viability score
\[
\mathcal{V}_{\mathcal{E}}(x) \in \mathbb{R}
\]
as any scalar functional such that:
\begin{enumerate}
\item $\mathcal{V}_{\mathcal{E}}(x)>0$ implies $x\in A_{\mathcal{E}}$ and all stated
      inequality contracts of $\mathcal{E}$ hold at $x$;
\item $\mathcal{V}_{\mathcal{E}}(x)\le 0$ implies at least one contract of $\mathcal{E}$ fails at $x$.
\end{enumerate}
\end{definition}

\begin{remark}
$\mathcal{V}_{\mathcal{E}}$ is not unique; it is a diagnostic that packages the
inequality contracts of $\mathcal{E}$ into a single scalar.  Scores for
different encodings are comparable only after a common normalization is
declared.
\end{remark}

\begin{definition}[Maximal viability]
Define the maximal viability at $x$ by
\[
\mathcal{V}^\ast(x) := \sup_{\mathcal{E}\in \mathrm{Enc}(x)} \mathcal{V}_{\mathcal{E}}(x),
\]
with the convention $\mathcal{V}^\ast(x)=-\infty$ if $\mathrm{Enc}(x)=\varnothing$.
\end{definition}

\begin{remark}
With a common normalization, $\mathcal{V}^\ast(x)$ is a best-margin
diagnostic.  Relating its sign or continuity to
$\partial X_{\mathrm{enc}}$ requires additional regularity of the encoding
family.
\end{remark}

\subsection{The overlap nerve of the atlas}

We now define a global topological summary of chart overlap.

\begin{definition}[Overlap relation]
Two encodings $\mathcal{E}_i=(A_i,\dots)$ and $\mathcal{E}_j=(A_j,\dots)$ overlap if
\[
A_i\cap A_j \neq \varnothing
\]
and there exists an exact re-encoding morphism or a controlled arrow with an
explicit error budget on the overlap.
\end{definition}

\begin{definition}[Atlas nerve]
Let $\{\mathcal{E}_\alpha\}$ be a chosen family of encodings.
The nerve $\mathcal{N}$ is the simplicial complex whose:
\begin{itemize}
\item vertices are encodings $\mathcal{E}_\alpha$;
\item a $k$-simplex $(\mathcal{E}_{\alpha_0},\dots,\mathcal{E}_{\alpha_k})$ exists
      iff $\bigcap_{i=0}^k A_{\alpha_i}\neq\varnothing$ and transition morphisms are
      controlled on that intersection.
\end{itemize}
\end{definition}

\begin{remark}
The nerve records the combinatorics of a chosen cover.  A hole in the nerve
is not automatically a gluing obstruction; that conclusion requires the
appropriate good-cover hypotheses and a nontrivial transition cocycle or
cohomology class.
\end{remark}

\subsection{Universal reduced relation}

Independently of whether deterministic descent holds, the typed data define a
global relation.

\begin{definition}[Universal reduced relation]
Define
\[
\mathcal{R}_t:=
\{(P_0x,P_t\Phi_tx):x\in X_0\}\subseteq Y_0\times Y_t.
\]
\end{definition}

\begin{remark}
$\mathcal{R}_t$ is always well-defined.  It is the graph of a single-valued
map $F_t:Y_0\to Y_t$ exactly when the descent
criterion~\eqref{eq:descent} holds.  A representative section chooses one
value from compatible upper representatives; it does not make the full
relation a graph when descent fails.
\end{remark}

\subsection{Total-atlas summary}

The total atlas consists of:
\begin{itemize}
\item the total encoding projection $\pi:\mathbb{E}\to X_0$;
\item the normalized viability diagnostic $\mathcal{V}^\ast$, when a common
normalization is supplied;
\item the overlap nerve $\mathcal{N}$ recording the chosen cover and its
controlled transition data;
\item the typed relation $\mathcal{R}_t\subseteq Y_0\times Y_t$.
\end{itemize}

No single global admissible encoding is asserted.  The chart-of-charts
records local availability without converting local contracts into a global
physical law.


% =====================================================
\section*{Conclusion}
\addcontentsline{toc}{section}{Conclusion}

\noindent
Modal Triplet Theory is presented here as a conditional, non-self-sealing
framework. It asserts no ontology and no universal effective law.
Instead, it gives checkable gates for representative selection, exact
recovery, autonomous descent, basin-local stability, controlled
re-encoding, and measure-dependent probability.  Finite admissibility is a
stated axiom about certified encodings, not a consequence of noninjectivity.
Every stronger conclusion must identify the margin, measure, dynamics, or
physical source theorem that supplies it.

\vspace{1em}

\noindent
Any physical interpretation is an external assignment layered atop the
mathematical structure developed herein.
\bigskip
\nocite{Nero2026_MTT_Admissibility_Encodings_Structure,Nero2025_MTT_Foundation,Nero2025_FixedPoints_I,Nero2025_FixedPoints_II,Nero2025_FixedPoints_III,Nero2025_FixedPoints_IV,Nero2025_FixedPoints_V,Nero2025_FixedPoints_VI,Nero2026_ProjectionAdmissibility_Principle}
\bibliographystyle{plain} % or abbrv, unsrt, alpha, etc
\bibliography{main}

\end{document}
